%% =====================================================================
%%  B4-T-15-01 -- what B4 contributes to "The Crowd as Evidence"
%%  -------------------------------------------------------------------
%%  LAYER A: APPEND-ONLY. Written once, then left alone. Material found
%%  only here sits in the `unique` slot and is protected by rule R10.
%% =====================================================================
%% @kind: source
%% @id: B4-T-15-01
%% @book: B4
%% @topic: T-15-01
%% @locator: ch. 4, pp. 126--156
%% @status: distilled
%% @unique: yes
%% @integrated: yes
%% @updated: 2026-09-19

\dossier{B4}{T-15-01}{\bookshort{B4}}{ch. 4, pp. 126--156}

\begin{distilled}
  \item What many people are doing is taken as evidence of what is correct, and the inference is followed hardest where personal expertise is least.
  \item The behaviour of a crowd is treated as independent sampling when it is usually mutual copying.
  \item Numbers can be displayed rather than produced, and the display carries most of the persuasive weight.
  \item Applied to audiences, to queues, to testimonials and to the presentation of popularity as a fact about quality.
\end{distilled}

\begin{unique}
  \item The independence failure: the count does not multiply the evidence because the members are watching each other. This is the analytical core of the entry and is not present in B1 or B3.
  \item The resulting asymmetry --- the number need only be displayed, not earned.
\end{unique}
